\section{Introduction}

Every day, millions of transit travelers consult routing
applications---Google Maps, Naver Maps, or open-source engines such as
OpenTripPlanner (OTP)---to decide how to get from origin to destination. These
platforms uniformly adopt a deterministic paradigm: they minimize a predefined
generalized cost and present the result as the single ``best'' route,
implicitly assuming that all users weigh walking, riding, and transferring in
the same way. Empirical evidence, however, paints a different picture. Elderly
passengers may tolerate longer rides to avoid stairway transfers; wheelchair
users may detour to reach an elevator; commuters may prioritize speed above all
else. If these preferences are systematic rather than random, a
one-size-fits-all recommendation will systematically mis-serve certain
groups---a concern with direct equity implications as urban populations age.

Seoul provides an ideal setting for investigating these questions. Its
Transportation Card system records over seven million tap-in/tap-out
transactions daily across all bus and subway services, enabling
population-scale revealed-preference observation of route choice. Five
administratively defined card types---general, elderly, disabled, youth, and
children---permit demographic segmentation without survey recruitment, and
OpenTripPlanner's RAPTOR algorithm \cite{delling2015} can systematically
generate competing alternatives for each observed trip. Prior work using Seoul
smart card data \cite{kim2020} calibrated a transit route choice model but
considered only a single user type; the present study extends this to four
model families, five user types, and 1.32 million trip chains.

Despite a growing literature on transit route choice, three substantive gaps
remain. First, most studies estimate a single discrete choice specification and
report its fit, without testing whether the structure imposed by random utility
theory helps or hinders prediction relative to flexible machine learning
alternatives such as gradient boosting. Second, neither Mixed Logit's
continuous coefficient distributions nor observed demographic segmentation can
discover data-driven behavioral clusters that cut across administrative
card-type categories---a 70-year-old commuter who walks briskly differs
fundamentally from a frail peer who avoids stairs, yet both carry the same
``elderly'' smart card. Third, route choice research overwhelmingly concludes
at the point of parameter estimation; whether those parameters can improve an
operational routing system---and if so, by how much---is seldom examined.

The present study pursues four objectives. First, we develop a multi-level
route similarity framework with 10 metrics organized into four hierarchical
levels. Second, we estimate and compare four model families---MNL, Mixed
Logit, Latent Class (with concomitant variables), and LightGBM---on an
identical hold-out test set drawn from 1.32 million trip chains. Third, we
characterize the latent behavioral segments and examine how they align with
administrative card-type categories. Fourth, we close the loop between
estimation and practice by demonstrating that MNL-based re-ranking of
router-generated alternatives outperforms traditional parameter calibration at
negligible computational cost.
